\subsection{На сколько применимы \acs{ed} для \acs{la}?}
Ключевые моменты \{ Дроны, Гражданская авиация, экология, простота 
конструкции и простое обслуживание по отношению к ТРД, высокий кпд\} \\

Электродвигатели в авиации: будущее уже здесь? \\
Современная авиация стоит на пороге революции: электродвигатели постепенно вытесняют традиционные турбореактивные и поршневые силовые установки.
Основные их преимущества --- это высокая экологичность, простота конструкции и обслуживания, а также высокий КПД.
В 2025 году лидерами электрификации стали \acs{bpla}. В Силу расширения электронной комплектной базы беспилотники уже массово используют нефтегазовые компании, 
ученые экспедициях, в военной отрасли.
А вот пассажирские самолёты на электричестве всё еще редкость, существуют гибридные системы и полностью электрические модели которые уже испытываются. 
Например, проекты вроде Alice от Eviation показывают, что региональные перелёты на батареях возможны. 
Но комерческая реализация магистральных самолетов дальнего и среднего класса на сегодняшний день затруднена(подробнее см. в разделе про литий-ионные батареи)

В чем же преимущества \acs{ed}?
Электродвигатели не выбрасывают ${CO}_2$, имеют низкое шумовое загрязнение. Их обслуживание проще: 
    меньше подвижных частей, 
    нет необходимости в сложной топливной системе и масле. 
    Это снижает эксплуатационные расходы, что особенно важно для коммерческого использования.
\subsection{Различные источники ЭДС}

\subsubsection{Литий-ионные батареи}
Ключевые моменты \{Удельная мощность, токоотдача(45C), популярность \}
\cite{schafer2019technological} - экономически выгодная плотность энергии ячеек.
https://www.nature.com/articles/s41560-018-0294-x
\subsubsection{Водородные топливные ячейки}
Ключевые моменты \{ большая удельная мощьность по сравнению с li-ion, сущевствование комерческих решений,
аналогичные методы контроля элементов \}